\chapter*{下编\quad 考释}
\phantomsection\addcontentsline{toc}{chapter}{下编\quad 考释}
\markboth{下编\quad 考释}{下编\quad 考释}
\chaptersubtitle{ЧАСТЬ ВТОРАЯ}{ИССЛЕДОВАНИЕ}
\setcounter{chapter}{0}

本编凡八章。每章均自上编所收之一片或数片残简立论，旨在厘清该条材料所能证实之范围与所不能证实之边界；凡上编校记已经解决之文字问题，本编不再重复，而上编留存之史料空白，本编亦不强加弥合。八章之间实非平行罗列，而系依推演次序逐层展开：第一章论居民之体质与制度定位，第二章论其居所之形制与尺度，第三章论维系此城运转之净秽制度，第四章专考该制度中之一项特例，第五章论其核心收集机构与劳役群体，第六章汇总岁计与人口账目，第七章还原撤离之经过，第八章则归结于名籍关系。

本编自始至终严守一道界限：凡三十七片残简所未载者，一概不予纳入。传世文献中与月面城邑牵连之诸多叙事与名号，本编均不涉猎，非因其毫无研究价值，乃因本书体例既定为辑佚与考释，考释之对象自当严格限定于所辑出之三十七片实物本身。阅读次序上，第二章与第七章可先行检阅，前者据实测尺度复原城池之空间格局，后者则串联撤离前最后一年之档案序列。

\chapter{月之民的体质与制度}

简十四将「计口三十一万二千」与「实居十一万九千」两项数字并列刻于同一残片之上。就统计口径而言，两数所指本非同一时段，然刻写者将其紧接排布，恰足以表明在当时之行政意识中，二者必须互相参照方能读通。两数之间相距十九万三千之差额，其中十六万二千人系建城期之减员：依莫斯科档案馆第 117 全宗第 2187 号卷宗所作先期复原，该城开工规划之容量为三十一万二千人，迨至竣工之时，实存人数已降至十五万；余下之三万一千人，则为竣工之后至公元前二百二十一年撤离为止的漫长时段内陆续减少之数。刻写者将额定与实居并列而不加修饰，正说明其熟知该项落差非但存在，且早已视为常态。

\section{来源}

关于月之民之具体来历，三十七片残简本身并无任何直接记载，唯一可以确证者，仅在于其抵达月球正面静海地区之时间远早于这批残简刻成之前，且其初至之时，当地尚处于未经开发的原始状态。此项判断的依据见于简一所记：「以石为垣，垣厚四丈，高十二丈」，又记「地无草木，风自西来，终岁不息」。此数语所记录者乃一处营建伊始之基地环境：墙垣之具体厚高、风向之恒定方向乃至地表草木之完全缺如，皆须刻于金属留档，原因仅在于其时一切皆属初创。生活在故土之人群断无必要将常年风向与地表植被状态刻入行政简牍之中，惟有面对新辟异域之初创者方会以此为建档之起点。简一所记乃营建而非迁徙本身，由此观之，当这三十七片档案开始刻写之际，最初之迁徙历程早已成为遥远之前事。

\section{体质改造}

简十三规定：「新入者，先居外舍三月，乃入环。入环之日，给籍，给衣，给器。」此条所指「新入者」，系指自地球表面往返归来之人。凡入环必先于垣外隔离居留满三月之期，其性质显非普通门禁查验所能解释；倘若仅仅跨越一道城门，断无必要设置长达九十日之缓冲滞留。此三个月内究系施行何种处置，残简并无一字明言，然参以简四与简五所载气路、水路、外工及浴水处置诸条（「出则浴，浴水不入渠」「衣别贮，不与内同」），可见凡自环外往返之人，皆须经历换气、换水、专渠洗浴及衣物分储等一整套繁复程序，继之以三月外舍居留，方准进入居住环。此等隔离规程之严密程度，与寻常尘埃防护之需求全然不成比例，其所防范者必为更深层之体质与环境冲突。

\section{改造的代价}

月之民在入城前究竟实施过何种自身改造，本书所辑材料仅保留了结果而完全隐去了工程细节，因此仅能由其制度后果加以反推。据相关档案推知，该套改造系一次性完成，其核心技术在月面设施竣工前不久即已失传，此后月之民亦不复能够增殖同类。此项改造之直接后果，在于使其生理状态脱离自然变化之序列：不再老化，不再具备生育能力，亦不再对新环境产生演化意义上的适应；作为交换，其个体寿命得以延长至足以维系同一座城池运转两万年之久。需要指出的是，此种改造与入环登记原非一事。简十三「给籍，给衣，给器」之规定，仅表明入环是一套标准化的行政手续而非惩罚机制，并不存在强制囚禁之痕迹。至于简十五所载「不足则减，减自上始」，其削减次序为先居者、次工、最后匠；然而同一残片所记口粮定额则为「人日二升，匠倍之，外工又倍之」，即口粮最低者为居者，最高者反为外工。配给序列与削减序列彼此相背，恰足以说明所谓「上」，在该城之行政秩序中并非一个固定不变之特权位阶。

\section{寿限}

城中居民之寿命界限，见于简十九所录「兔乱二万二千」之实亡与简二十六「出界则复老」之界说，与莫斯科档案馆第 117 全宗第 2187 号卷宗第 96 叶所录地面四代留存率曲线分属两事，不可混为一谈。第 96 叶属于后来归并之地面专题统计，本卷三十七片原件之中并无一字涉及落地后四代之数；相反，简二十六明言「凡去者，出界则复老，不复入环，亦不复通音书。已去之籍，唯存其数，不记其后」。「出界则复老」一语，足以从月面原件自身证实现世所谓不老之态纯粹受制于城内界限以内之环境，一旦离环赴地，便即重新归入凡人之衰亡序列；而「亦不复通音书」「不记其后」八字，更直截斩断了撤离之后月面犹然续记地面人丁之可能：名籍至此唯留总口数，自此以后，月面原档于地面后裔概无一字留存。

\section{器物上的旧主}

在考察月之民制度的整体冷酷性时，简十三所留存之一处细节颇与全书气氛相悖：「器不继，以旧器补。旧器有记，记其旧主。」简二十二第十七项明言「名，削」，足见该体系以消解个人名籍为常态；然而在器物配给之日常实务中，旧器之上却偏偏保留原主之标记。笔者并不认为此处构成逻辑矛盾：撤离时削名属于体制终结之际的清算，日常中记名则属于物资流转之追溯惯例，二者各司其职；但两种截然相反之原则同时并存于同一行政体系之内，仍足以说明该制度对「名」的态度并未达到完全自洽的程度。

较器物记名更为罕见者，乃是简五末段所记一桩具象个案：一名外工因越界被捕，审问其缘故，仅答「欲见山」三字。在三十七片悉数由规程与账目构成的档案之中，此三字是唯一一句出自普通役夫之口的话语，其人职名与姓氏皆无存留，唯此三字缘由传世。按简五所记，该名外工平日居于垣外小舍，每月入环不得超过三个时辰，既非掌册之吏，亦非在编之匠，却宁可冒越界之刑（按简九规定，越界在此体系内与直呼山名同罪）而往观垣外之山。此三条简文之历史价值与文学抒情毫无关系，其意义仅在于证实：即便在一座以消除变化为最高原则而设计出来的城池之中，制度仍未完全根除个体向外张望之冲动。关于此条涉及名与界的法律关联，本书将在第八章再作深究。

\chapter{城的形制与尺度}

前章既已考订居民之体质与身份定位，本章继而复原其栖居空间之物理形制与工程尺度。城池平面规模之唯一直接线索见于简六「垣周九里」四字；按秦汉尺度一里约合四百一十六米折算，\footnote{秦汉一里三百步，一步六尺，一尺二十三点一厘米，实合四百一十五点八米；此处取四百一十六米，旨在便于读者以整数复核推算，若改按四百一十五点八米计，全周长相差不足五米，于量级判断无涉。}九里之周长约合三千七百四十四米，由此反推其直径为一千一百九十二米，半径五百九十六米，占地面积一点一二平方公里。简三明言城内分作六环，若将五百九十六米半径均分为六段，则每环径向宽度约为九十九米。此石垣所围合者仅系核心居住区，城之外缘尚设有一道直径约三点五公里之屏蔽环，垣与屏蔽环之间为空旷缓冲地带，仅近垣处设有外工居舍（详《月面设施复原图集》图版 11 与第 117 全宗第 2187 号卷宗第 12 叶总平面图）。本书以下所论各数，除另有注明者外，皆以主垣以内为边界；这些尺度实非凭空旁证而来，而是将残简所记古制单位换算为现代公制之后，各条零散记述自然显露出的空间秩序。

\section{六环}

简三所列六环实非六个彼此平行的独立街区，而是同心嵌套之环状结构。简二称「三环以上，居者稀」，由此可知环次之编号方向系由外向内递增：一环居最外缘，六环居最核心。居住人口主要集中于外侧一至三环，每环宽九十九米，三层合计纵深约三百米；三环以内居者渐少，至于最内层之六环，则当为整座城池之中枢所在。若将撤离前之十一万九千实居人口摊分入此一点一二平方公里之内，其平均人口密度高达每平方公里十万六千六百人；即以当代地球人口最稠密之都市计，最高密度亦不过每平方公里四万六千人上下，\footnote{此数取近二十年各国都市人口统计中之最高档次；鉴于各国城市乃至同一城市内部各街区之间密度差异甚巨，此处仅作量级比照。}该城之密度竟达其二点五倍。尤须注意者，此十万六千余人之密度尚属该城「空置六成二」之后的最宽松状态：按简十四，城池额定容量为三十一万二千人，而两万年间人口只减不增，撤离前夕乃其历史上人数最少之时；即便如此，按十一万九千人平均分摊，每位居民所占有之平面面积亦仅九点四平方米，且此九点四平方米尚须涵盖寝卧、饮食、通道及一切公共设施之公摊份额，其空间压迫程度之反常，不待繁言而自明。

\section{门禁}

简一载「环各有门，门不可同启」，简三载「环各一门，门各有籍」。六道门户各有独立名册，彼此互异，且严禁同时开启，这一制度直接将城内各环之间的交通速率限制至极低水平。居民无法由一环径直行至六环，每越一环皆须先出本环之门、再入次环之门，且每次过门均须经过一次名籍核验。简二所谓「一环之口，常不过千」，即指此类过门查验口而言，谓一道门户同一时间所容留核验者常不足千人。姑以均值每环一万九千八百人计（实则三环以上居者稀少，外三环均值远高于此），即便假定一道门每小时可通过千人，全环人口依次通过亦须耗费二十小时。由此可见，这套门禁系统在设计之初便非用于维持日常通勤，其真正职能在于严格管控跨环进出；至于维持城池运转所需的日常物料与工程往来，势必另有隐蔽通道独立运作。

\section{气路与水路}

关于环与环之间的暗道运行方式，简四提供了唯一一段工程记述：「气自下而上，水自上而下，皆以管。管在壁中，不可见。」气体自下上升，水流自上而下，皆密封于隔墙管壁之内，这一表述表明城池内部在竖直方向上存在着明确的高程差与压差结构，即气源设于底部而向上升腾，储水设于顶部而向下输送，绝非平铺展开之六层平房所能容纳。一座直径仅一点二公里之圆城究竟凭借何种机械或物理机制维系此种垂直循环，残简并无一字说明，唯一可以断定者，在于它是一个具备封闭竖直循环系统的工程整体。在全书所有工程细节之中，「管在壁中，不可见」七字分量尤重：居住在环内的人无法目睹维系自身呼吸与饮水之管网，也就无从直接感知自己赖以存活的物质基础，整座城池的生命维持体系对其居民而言是彻底隐形的。

\section{日常}

除空间形制外，城中居民每日之作息节律亦无任何完整日记体残简留存，唯一可行之路径，乃是将散见各片之动作条目依时段拼合复原。清晨之际，六道环门中之一道开启，三百二十名外工由垣外小舍入内，依规定每环仅通行单一通道，且滞留不得逾三个时辰（简五）；其具体作业内容残简未明言，参以简十三「给器」与简六「限外不得置物」诸条推之，当以巡修为主。日间之时，三千名「捣者」在环外某处特定工坊内依每日三合之定额劳作（简十六、简十七），达额者赏粟，不足者罚役，连犯三日以上者则改移名籍，此即全书唯一能窥见常态劳动节奏之条文。

至于周年岁计层面，一年之中仅有两项固定公事：其一为气路之年度检修，「岁一检」，检者不得入环，仅于环外通道作业，检毕将读数录于木板并悬诸门口（简四）；其二为名籍之年度更替，「岁一易」，旧籍于固定日期由二人焚毁、三人监视并五人联署签名（简十八）。除上述开门、做工、检修、焚籍四类公事之外，三十七片内绝无任何节庆、婚丧或日常交际之记载，唯一流露个体声音者，仅有前述越界外工被捕时所答「欲见山」三字。笔者并不主张凭此断定城中绝无任何多余生活（残简之本质乃是账册与规程而非日记，账册本不记节庆，规程亦不录闲谈），但笔者同样拒绝以「想来应当亦有」之揣测来填补史料之空白。就文献呈示之实况而言，在一个以排除变化为建城宗旨的封闭体制内，时间之流逝势必高度均质化：既留驻了恒定之温候与无别之昼夜，便已将变化本身剔除在日常之外，三十七片之所以读来千篇一律，恐怕并不单纯因为它们在文体上属于账簿。

\section{光与声}

在所有未曾直接见诸文字的日常要素中，唯「光」与「声」两项尚可由既有残简反向推出，且此二项最足以界定城内之感官环境。就光而言，简一明确记「昼夜如一」，且全书三十七片纪时从不用「日」与「时辰」（唯简五外工入环限「三时」为例外），仅以「岁」「月」纪事；两证相合，只能得出同一结论，即昼夜交替之物理节律在环内并无实质功能，至于是依靠某种人造光源维持长明，抑或处于某种恒定微光乃至长暗之中，残简无一语可凭。就声而言，简一所写「风自西来，终岁不息」专指垣外月面或屏蔽层间环境，而垣内则是六环重垣相套、六门永不同启、气水管道深埋壁内；此数项条件叠加，其所构成的必然是一座近乎绝对静寂之所在。三十七片内绝无一字记及声音，全帙唯一含带听觉反衬的记载，反而在撤离终结之简三十一：「门闭之日，无送者，无哭声」，那已是城池生命终结的最后一日。光与声在档案中的同时缺席原非史料遗漏，盖在本来就极少记事的账册体系里，这两项恒定不变之背景条件恰恰是最无须入账之物。

\section{垣外}

自城心向外部推演，其空间序列依次为：六环中枢、内三环、外三环居住区、石垣、垣外百步空地（依简六严禁置物、起舍、安葬）、外工小舍（容纳三百二十人，见简五）、以及「界」（简九），界外则复有「山」峙立。由中心至「界」之全长尚不足两公里，其中一千一百九十二米为城池本体，至于界外至山之距究为几许，残简并无一言。人若自最内侧之六环徒步出至界外，在空间尺度上无异于将整座城池彻底穿行一遍。然而简九偏偏记称：「在其内者出，必失其半。失者不复。」倘若此句确属实写而非修辞，则意味着穿行这不足两公里之距离并跨越那道「界」，必将付出某种不可逆之半数损耗；至于所失之「半」究为何物，残简同样不着一字。

\section{本章方法的说明}

行文至此，有必要就本章推算方法之限度作一明确说明。本章所复原之城池几何尺度，全数推演自简六「垣周九里」四字，其公开算式取秦汉一里约合四百一十六米、圆周率取三点一四一六。读者尽可自行验算或改用其他里制常数：倘若将一里径取作五百米，则全城直径相应增至一千四百三十二米，人口密度随之降至每平方公里七万四千人，全套数字均会随这一笔者无法最终核实的古制常数而在百分之二十幅度内上下摆动。尽管如此，本章之核心判断并不依附于这一常数而动摇：无论一里取四百一十六米抑或五百米，此城在尺度上之狭小可穿、在人口密度上之违背常理、以及其维生系统对居者而言之彻底隐形，此三条结构性结论均完全成立。

\chapter{净秽制度}

简五所记「出则浴，浴水不入渠」十字，若抽离全书语境而孤立视之，极易被误读为一条寻常的卫生防疫规定；然将其置于三十七片之整体脉络中加以审视，此条实为揭示整座城池运行逻辑之入口所在。本章以下将由气水管道之工程维护切入，继而辨析环内「洁净」与外部「污秽」之政制界定，试以三十七片所存实证条目取代后世传闻之先验预设。

\section{洁净的维持}

月之都之洁净绝非天然禀赋。这一判断必须先行确立，盖古往今来几乎一切关于月面圣域之传说，皆先验地假定其本自无垢。然而简四所记「气自下而上，水自上而下，皆以管。管在壁中，不可见。凡管，岁一检」，清晰地勾勒出一套必须依赖工程维护方可运转的人工循环体系：天然无垢之地无须管网，更无须以「岁一检」的制度化周期予以强制监控。不仅如此，简四紧接记明：「检者不入环，于道中为之。检毕，录其数于板，板悬于门。」检修人员严禁步入居住环，仅得于环外通道作业，且完工后必须将检测读数公开悬榜于门首。若仅为常规管道修缮，绝无将技术数据公之于众的必要；将工程检测上升为公开档案记录，正足证明维系环内洁净乃是该城政制之首要基石。

\section{污秽的定义}

三十七片之中并无任何一条直接为「污秽」二字下定义：简五仅载其隔离与洗浴之处置，简十九仅载其引发之事故死亡数，简二十二第十一项则仅余「山，封」二字。尽管缺乏抽象界定，所有相关条文却均严密指向同一个核心目标，即把一切会产生「变化」之要素阻隔于城外，或将已被带入城内的变化要素彻底清除出境。空气必须经由管道过滤，水源必须封闭循环，外工出环必先洗浴且衣别贮藏，垣外百步禁绝置物、起舍与安葬，乃至年度换册焚余之灰亦须远埋垣外。在这一整套严苛得近乎偏执的规制面前，物理意义上的灰尘与微生物至多只是表象；在该体系的内在逻辑里，「脏」与「变」实属同一概念，「净」与「不变」亦属同一概念。遍检三十七片，笔者未能发现任何足以推翻此项等价关系的反例。

\section{物与人}

既然在此体系内「污秽」即等同于「变化」，其防御制度便必然围绕「物」与「人」两类载体展开。关于「物」之处置，最典型者见于简六与简十八。简十八记「凡籍，岁一易。易籍之日，旧籍焚……焚余之灰，埋于垣外」：旧籍纵然已被焚为灰烬，仍须逐出垣外深埋，足见制度所要排除的并非纸片或金属载体本身，而是载体之上所附着的「旧」与「过去」这一时间属性；将旧物彻底清理出城，正是岁焚制度的根本目的。关于「人」之处置，则以简五与简十三最为显豁：外工月一入环、出则浴且浴水不入渠、衣别贮；新入者须先于外舍居留三月方准入环。两套手续程序虽异，所针对者却是同一类对象，即刚自外部变化世界返抵之人，目的皆在于将其携带之「变化」剥离或隔离于居住环之外。

此处必须附带作出一项界定：本编在讨论净秽制度时，刻意不将「言」计入物与人两类之内。简九禁「呼山名」一事，初观之仿佛同属隔离禁忌，实则所涉乃名籍管理权限，与本章所论变与不变之物理净秽体系实属另一套机制；二者若强行混谈，必致两边逻辑俱乱，故简九专归第八章处理。

\section{界}

简九实为全书最难索解之一片，其文曰：「界者，非垣非墙，非水非火。在其内者出，必失其半。失者不复。」究竟何谓「界」，残简并无一字正面阐明：前四句纯用否定式排除法，仅告知读者其既非石垣、土墙，亦非水火阻隔，而否定式本身从不提供正面定义。「必失其半」所指为何同样无从判定，或指随身携带之物，或指跨越者自身之某种组成部分，唯有末句「失者不复」确认了该减损具有不可逆性。面对此类彻底沉默的档案，强行构拟解释并无学术意义；这一空白就悬置于此，成为比简二十六「唯存其数，不记其后」更为彻底的一处断裂，盖后者尚且指明了所缺为何事，而此处连所失为何物亦未尝留下只字片语。

\chapter{药，弃}

简二十一至简二十二所录退役清单共计十九项，各条均配以明确之工程或行政处置动词。本章仅聚焦于其中最为孤立的一条，即简二十二所载第十三项：「药，弃」。全书三十七片关于此物之直接记载，尽在于此二字。

\section{一个「弃」字}

在十九项退役举措中，「弃」字仅此一见。其余十八项之处置动作大致可归入封闭（环、库、山）、切断（气路、水路、界）、焚毁或损毁（籍、器）、熄灭（炉）、平夷堵塞（外舍、垣门）、收缴移除（册、界石）与削减消除（标、名、口）等数类，此类动作皆在既有行政制度之常规权限内运作，各有明确的执行者与验收标准，唯独「弃」游离于全部既定分类之外。在词义与工程实务上，「弃」与「毁」必须严格划清界限：毁意味着彻底破坏使其丧失原有功能，弃则意味着将其遗留原地、放弃管辖而不加销毁，被弃之物在撤离后仍存留于某处，只是不再随队伍携行。一个严谨到连旧籍灰烬都要深埋垣外的行政体系，在最后撤离时对某种物料下达「弃」的指令，其背后无非出于三种可能：无法携走、无需携走，或不愿使其留在自己手中；此三者亦可同时并存。此外，第十三项紧随第十二项「炉，熄」之后、第十四项「册，收」之前，中间仅与第十一项「山，封」相隔一项。清单序列通常反映现场处置之先后次序，倘此序可信，则对「药」的遗弃恰发生于炉火熄灭之后；也就是说，当执行遗弃之时，生产此物之核心装置已然停止运转。

\section{药是什么}

至于「药」究系何物，残简既未署其专名，亦未载其原料产地、制备配方、具体效用与服用方式；在全部三十七片之中，「药」字只在该项清单内出现过这一次。唯一可供比对的间接线索见于卷二：简十六记工籍「捣者三千」，简十七继记「定额：人日三合，多者赏，少者罚」。「捣」是一道常设工序，既有明确人数编制，又有严格之产量定额，且持续运行达两万年之久；其成品所用之计量单位为「合」，而汉代容量单位「合」专用于计量粉末或细小颗粒，断不用于计算器物件数。由此推论，倘若「捣」这道持续两万年的工序所产出者正是「药」，则整座月面城池在漫长存续期内唯一真正投入规模化「生产」的产品，便非此物莫属。当然，此说在文献上尚属孤证推演：三十七片中没有任何一片将「捣」与「药」直接相连，两个语词从未在同一残片内共现；笔者之所以仍提出此项关联，仅因遍检全书实物账目，再无第二个可与之匹配的候选物。

\section{一座只生产一样东西的城}

若将三十七片所见全部物资清单逐一排比，整座城池便呈现出一种颇为反常的经济面貌：所有档案皆指向输入、分配与消耗，粟、水、气三类储备由外部或地下系统维持，器物不足则以旧器补给，唯独找不到任何面向城外的产出，甚至找不到内部自给性制造的记录，唯一带有明确生产意味的记录便是三千名「捣者」日复一日地以「人日三合」之定额捣制某种颗粒或粉末。一座承载十一万九千居民运转两万年之久的庞大设施，却只为产出单一一项成品而设，且在最终撤离时将这项倾全城之力生产了两万年的唯一产物断然「弃」之于原地。这一悖论既可充作本章之结语，亦可视为悬而未决之疑案；笔者虽倾向于认作上述生产—遗弃链条的成立，然就现有三十七片材料而言，并无第二片残简可供交叉印证。

\section{本章的方法限度}

全章所考不过简文二字，围绕此二字展开四节推论，其中半数以上属于间接推演而非直接实证。笔者本可将该项径直并入第七章十九项退役清单中一带而过，之所以不避单薄而强行独立成章，仅因其牵涉全书一个更根本的结构性疑题：倘若一座城池两万年仅专注于生产唯一一种物质，却在终局将其整批抛弃，那么该物质究竟具有何种性质，便绝非可以轻描淡写略过的细枝末节，尽管仅凭目前这三十七片残简，本书尚无力回答此问。

\chapter{净秽山与月兔}

本章虽在总纲中列为第五章，在考释方法上却与前后各章颇有殊异之处：其余各章率由文字直接推演制度，本章则须先按字面强行读解，再由数量与制度之间的硬性矛盾迫使读法转向改读。这一不得不进行的读法转换，正构成本章第二节之核心论证。

\section{位置}

就城池外围之地理要素而言，本书共给出垣、界、山三项。按卷一与卷三所见线索，居住区止于垣周九里（直径约一点二公里），屏蔽环在其外而残简未及；垣外百步禁置物、起舍、安葬（简六），再外为外工小舍（简五），再外为界（简九），界外即为山，山距界几何则残简无一字交代。由此可知，自城心至界全长尚不足两公里，即便加上界外至山之距，该山亦必去城不远，在空间上完全处于徒步可达之范围内。然而简九却明载「界外有山，山名不可呼」，既将其置于不可逾越之界外，复严禁直呼其名。一座近在咫尺且明确承担某种实际功能的「山」，何以甚至连名称都沦为禁忌，此乃本章首须面对之疑问。该疑问在单纯的地理或工程层面实无从作答，盖其根源不在地理边界，而在名籍账册：三十七片所记之名原非自然专名，而是账册条目，私呼山名无异于擅自为册外之物立条入籍，直接触犯吏员专掌名籍之权。此层法理联系将于第八章第一节详加辨析，此处姑且存疑。

\section{容量}

与位置之疑并行者，尚有一笔可作量级核算之账目。十一万九千居民在两万年间所产生并排出之一切代谢与变质物质，均须经由特定系统收集、压实并封存，其累计总量究竟几何，残简并无一处明言，唯一可供换算产量之工序唯「捣」一项。简十六记捣者三千人，简十七记定额为「人日三合」；按汉制一合约合二十毫升折算，三合即六十毫升，三千人每日所捣成品凡一百八十公升。若按两万年、每年三百六十五日连贯乘积计算，其累计体积约为一百三十一万五千立方米，约当边长一百一十米之正立方体，规模近似一座小型水库。此数诚然只能视作量级估算，盖其前提假定三千编制两万年从未缺额、定额从未更张、且人口递减全不影响产量，三项假定在历史上均难严丝合缝；但其至少划定了一个不容忽视的数量级边界，即该道工序在两万年尺度上所处理的物料总量必然处于百万立方米这一量级。

下面的改读即由此数量级矛盾而生。倘若该处果真是一座自然土石山丘，容纳百余万立方米之废弃物料尚不致直接冲突；然而一旦将其与制度条文对照，三重矛盾便立即浮现。其一，简二十二第十九项处置门户，第十一项处置「山」；门在垣内而山在界外，二者隔着简九所谓「在其内者出，必失其半」之界。倘若该山确系一处必须常年往返维护之独立设施，则维修工役势必频繁出界，而频繁出界与「必失其半」的不可逆代价绝无相容之可能。由是推之，要么此项出界代价仅在某种一次性状态转换中触发而不随每次往返重复发生，要么即属笔者对「山」字之位置与字义读解有误。即便勉强采信前一种缓和解释，仍无法回答第二个更棘手之疑问：一座容量达百万立方米级的独立巨型构筑物，何以在简二十二精细到连「界石，移」「标，除」都一一单列的退役清单中，仅仅占得「山，封」两个字，既不记封几处，亦不言如何封堵。

基于上述双重困境，必须提出第三种读法，亦即笔者认为在逻辑上最站得住脚的一种解读：此处所谓「山」，根本不是一座独立于城池体系之外的自然山体或孤立建筑，而是整套净秽收集系统在空间上位于外围的一段核心装置。从垣外仰观，该装置或许呈现为一片人为堆筑或地形隆起之坡面，其表层寸草不生而具备显著斜面，故在记录中被概称为「山」（即《Лунная прародина》所复原之「净秽山」）；而在实际工程功能上，它正是承担将城内污秽抽吸排出、压实并予以封存的终端收集段。简九之所以严禁呼其真名，亦随之冰释：一旦直呼其工程专名，便无异于向环内居民挑明此物之真实用途，而此恰恰是严禁居者知悉之机制机密。这一改读固然无法由残简本身获得直接证词，却具有压倒性的解释效力：凡将「山」作字面山岳理解所必然引出的三项制度悖论，在此读法之下皆消弭于无形。

\section{编制}

第三重矛盾系于人数编制。简十六记「捣者三千」，而简十四与简十九显示全城人口由三十一万二千剧减至十一万九千，总口数衰减逾六成，何以「捣者」之人数始终恒定为三千而未随人口比例缩减。依照前节之改读，此项矛盾亦随之消解：既然该设施属于维系整座城池恒常运转的系统收集段，那么服务于该系统之「捣」工序所对应的便是固定之设施机位与处理负荷，而非随居住人口浮动之消费需求；系统自身规模在两万年间未尝变更，则对应之在册编制数自然维持三千不变。同时必须指出，三千乃在册之法定编制数，未必等同于该工序在城外实际动用之全部劳动力：若全期果真只有三千人按每人日三合捣作，两万年累计总量仅百万立方米余，较之一座运转两万年之巨型封闭城池所需处理的总体规模仍显偏低，此间差额唯有待后文考察劳动群体真实身份时再作申论。另有一条颇为微弱的旁证见于简二十二：清单将「山，封」列于第十一项，位于「册，收」「名，削」之前而居全表之中段，排于居首之居住环与水气系统之后。退役顺序通常折射设施之重要性等级，然此一旁证论证效力颇弱，笔者姑且并列于此，只因现有档案中别无更佳佐证。

\section{处置方式}

统计简二十一至简二十二所见十九项退役处置措施，其中「封」占十项、「断」占四项、「毁」占三项、「弃」占一项、「移」占一项。山列第十一项，其处置动词恰为「封」。在工程处置语义中，「封」与「毁」存在根本分野：毁乃不可逆之彻底破坏，封则属于可逆性闭锁，凡封存之对象在技术上均保留日后重新启封运转之余地。由此观之，撤离者虽将器物摧毁、居住环封堵、外舍夷平、不复为整座城池保留再启之可能，却唯独为「山」留出了一条可逆的封存后路。此条判断尚可与简二十三末句参证：简二十三记垣西等三处石板覆土之上「草不活」，且守者称「此三处，冬日不冻」，此乃全书三十七片中唯一提及撤离封存之后地表仍存热异常之证据。即便垣西三处与界外之山未必指同一工程部位，两条记述合观，仍足以证实公元前二百二十一年撤离封存之后，该区域地下系统并未即刻彻底冷却，尚有某种处于休眠或低功耗维持状态之热源存留于地下。

\section{月兔}

关于「月兔」一称，必须先作文献边界之澄清：三十七片残简本身并无「月兔」连文，仅于简十六记「捣者三千」之工种名，于简十九记「兔乱二万二千」之事变名；「月兔」作为对该劳役群体的统称，见于巴兰尼科夫《Лунная прародина》（2018）一书对相关工籍残卷的跨档案比对，本书姑循此译名指代简文所见之「兔」与「捣者」。三十七片之中并无任何一条正面说明月兔之生物属性或社会身份，唯一之直接制度线索见于简十六：该片将其与「搬运」「修理」「杂役」并列归入「工」籍，并特别注明「捣者不入环，月给别」。不入居住环而口粮另行发放，这两项规定清晰界定了其在城内秩序中的位置：他们明确被纳入城池行政编制之内，却又被物理隔离于核心居住区之外，承担苦役且享有与环内居者截然不同的配给体系。

至于这一劳役群体的真实总人数，三十七片内绝无一字明载；若仅据简十六字面推之，极易误以为其总数不过数千人。然而简十九「事故之数：兔乱二万二千」一条，恰恰暴露了这一字面印象的虚妄。从人数逻辑言之，二万二千绝不可能是月兔自身在乱事中的死亡数，盖与「捣者三千」的在编人数绝不相容，盖一支仅三千人之编制断无可能在单次动乱中死亡两万二千人。因此，简十九所记录的二万二千死者，必然是指月之民在「兔乱」中的损失人数，占其撤离前十一万九千人口之近百分之十九。在一场规模惨烈到足以令城主群体损员近两成的暴动之后，档案精确记录了统治者一方的死亡数字，却对发动暴动之劳役群体自身阵亡几何、原有总数几何只字未提。笔者并不认为这是书吏记账时的偶然疏漏。简十六以「捣者」这一纯然工种名称替代其族群名号，本已将该群体从社会身份谱系中抹去；一个连名籍身份都遭技术性取消的群体，自然不会被纳入官方事故统计的分母之中。关于月兔群体之实际总规模，必须到本卷以外的其他档案中去追索，三十七片所能呈示的历史真相止步于此。

\section{八千之数}

在进入下一章汇总全城账目之前，必须预先厘清一笔极易混淆之数字关系：简十九记兔乱死者二万二千人，而按第 117 全宗第 2187 号卷宗所载完工时十五万人口减除六批合计迁离之十二万人，两万年间城内实减三万人；以三万减去二万二千，尚余八千人之差额。此八千人在全部三十七片内没有任何一条记载，既不列入事故统计，亦不在撤离名单或留守人口之内，下一章将专门剖析此项无法平账之缺口。

\section{本章的限度}

本章乃是下编八章中直接史料最显单薄的两章之一，其一手残简之分量本不足以独立支撑整章之篇幅：三十七片之中提及「山」或「兔」者合计不过六片，且文字皆颇为简短，若将其拆散并入前几章，读者在行文上亦不会感到任何缺憾。之所以仍勉强单列成章，盖因该二项对象在域外其他档案与后世文献中被反复提及，而在月面基地自身遗留的原始文书里却近乎隐没无踪。一座城池最核心之净秽终端，在城池自存档案里仅占六字；支撑其运转两万年之底层工役，全书亦仅留五字。这种极度悬殊的记录比例本身，即是一项比文字记载更为坦白的一手史料：它直接证明了城内档案之编纂原则，从来不按「何事重要」来决定取舍，而是严格依循「何人有权知晓」的层级界限加以筛选。
